% !TEX program = xelatex
% 공통수학2 · Ⅱ. 집합과 명제 · 1. 집합 · 03 교집합과 합집합 · 1/2차시
% 학생용 활동지 (답안 없음)
\documentclass[11pt,a4paper]{article}
\usepackage{kotex}
\usepackage{iftex}
\ifPDFTeX\else
  \setmainhangulfont{Noto Serif CJK KR}
  \setsanshangulfont{Noto Sans CJK KR}
\fi
\usepackage[margin=2.1cm]{geometry}
\usepackage{parskip}
\usepackage{enumitem}
\usepackage{array}
\usepackage{tabularx}
\usepackage{xcolor}
\usepackage{amssymb}
\usepackage{tikz}
\usepackage[most]{tcolorbox}
\usepackage{fancyhdr}
\usepackage{titlesec}

\definecolor{goldc}{HTML}{B07C22}
\definecolor{goldstrong}{HTML}{8A5F16}
\definecolor{indigoc}{HTML}{2B3B57}
\definecolor{indigosoft}{HTML}{6E82A6}
\definecolor{linec}{HTML}{D6DACB}
\definecolor{surface2}{HTML}{E4E8DC}
\definecolor{writeline}{HTML}{C7CCBA}
\definecolor{inksoft}{HTML}{52604F}

\titleformat{\section}{\normalfont\sffamily\bfseries\large\color{indigoc}}{\thesection}{0.6em}{}
\titlespacing{\section}{0pt}{14pt}{6pt}
\newtcolorbox{goalbox}{enhanced, breakable, colback=white, colframe=goldc, boxrule=0.5pt, leftrule=3pt, arc=2pt, left=10pt, right=10pt, top=8pt, bottom=8pt}
\newtcolorbox{sheetbox}{enhanced, breakable, colback=white, colframe=linec, boxrule=0.6pt, arc=3pt, left=11pt, right=11pt, top=9pt, bottom=9pt}
\newcommand{\writespace}[1]{%
  \par\nobreak\vspace{3pt}
  \foreach \n in {1,...,#1} {%
    \noindent\textcolor{writeline}{\rule{\linewidth}{0.4pt}}\par\vspace{0.62cm}%
  }
}
\newcommand{\fillblank}[1]{\makebox[#1]{\hrulefill}}
\pagestyle{fancy}\fancyhf{}\renewcommand{\headrulewidth}{0pt}
\fancyfoot[L]{\footnotesize\color{inksoft} 공통수학2 · 2026학년도 2학기 · 지평선고등학교}
\fancyfoot[R]{\footnotesize\color{inksoft} 다음 시간 --- 03 교집합과 합집합(2)}

\begin{document}

\noindent
\begin{minipage}[b]{0.62\linewidth}
  {\sffamily\footnotesize\color{indigosoft} 공통수학2 · Ⅱ. 집합과 명제 · 1. 집합}\\[2pt]
  {\sffamily\bfseries\Large 03 교집합과 합집합 \textperiodcentered\ 21차시}
\end{minipage}\hfill
\begin{minipage}[b]{0.36\linewidth}
  \raggedleft\small\color{inksoft}
  반\ \fillblank{1.6cm} \quad 번호\ \fillblank{1.6cm}\\[4pt]
  이름\ \fillblank{3.6cm}
\end{minipage}

\vspace{4pt}
\noindent{\color{goldc}\rule{\linewidth}{2.2pt}}
\vspace{10pt}

\begin{goalbox}
\textbf{오늘의 목표}\\
교집합과 합집합의 뜻을 알고 구할 수 있으며, 서로소를 판별할 수 있다.
\vspace{4pt}
\small\color{inksoft}
$\square$ 자신 있음 \quad $\square$ 조금 헷갈림 \quad $\square$ 처음 봄 \hfill (수업 전 나의 상태에 표시)
\end{goalbox}

\section{생각 열기 --- 요리 재료 두 집합}

\begin{sheetbox}
볶음밥 재료: 밥, 버섯, 당근, 양파, 소고기, 계란\\
계란말이 재료: 계란, 당근, 버섯, 파, 치즈

\begin{enumerate}[leftmargin=1.4em, itemsep=2pt]
  \item 두 음식에 공통으로 들어가는 재료를 모두 말해 보자.
  \item 두 음식을 만들기 위해 준비해야 하는 재료를 모두 말해 보자.
\end{enumerate}
\writespace{3}
\end{sheetbox}

\section{개념 정리 --- 교집합과 합집합}

\begin{sheetbox}
두 집합 $A$, $B$에 대하여

\begin{itemize}[leftmargin=1.6em, itemsep=4pt]
  \item $A$에도 속하고 $B$에도 속하는 모든 원소로 이루어진 집합을 \textbf{교집합}이라 하고, $A\cap B = \{x\mid x\in A\ \fillblank{2.4cm}\ x\in B\}$
  \item $A$에 속하거나 $B$에 속하는 모든 원소로 이루어진 집합을 \textbf{합집합}이라 하고, $A\cup B = \{x\mid x\in A\ \fillblank{2.4cm}\ x\in B\}$
\end{itemize}

\vspace{6pt}
두 집합 $A$, $B$에서 공통인 원소가 하나도 없을 때, 즉 $A\cap B=\ \fillblank{2cm}$ 일 때, $A$와 $B$는 \textbf{서로소}라고 한다.

\vspace{8pt}
\textbf{생각해보기} --- 모든 집합과 서로소인 집합은 무엇일까?
\writespace{1}
\end{sheetbox}

\section{연습 문제}

\begin{sheetbox}
\textbf{1.} 다음 두 집합 $A$, $B$에 대하여 $A\cap B$와 $A\cup B$를 구하시오.\\
\hspace*{1.6em}⑴ $A=\{1,3,9,27\}$, $B=\{3,9,15,21,27\}$\\
\hspace*{1.6em}⑵ $A=\{x\mid x^2-3x-4=0\}$, $B=\{x\mid x$는 5 이하의 자연수$\}$
\writespace{5}

\vspace{6pt}
\textbf{2.} 다음 두 집합 $A$, $B$가 서로소인지 아닌지 확인하시오.\\
\hspace*{1.6em}⑴ $A=\{x\mid |x|\ge1\}$, $B=\{x\mid x^2=1\}$\\
\hspace*{1.6em}⑵ $A=\{x\mid x$는 정삼각형$\}$, $B=\{x\mid x$는 직각삼각형$\}$
\writespace{4}
\end{sheetbox}

\section{사고의 가시화 --- See · Think · Wonder}

\noindent{\footnotesize\color{inksoft} 오늘 배운 `교집합과 합집합'을 떠올리며 아래 세 칸을 채워 보자.}\\[6pt]
\begin{tabularx}{\linewidth}{@{}XXX@{}}
\textbf{\footnotesize\color{indigoc} See --- 무엇을 보았나?} &
\textbf{\footnotesize\color{indigoc} Think --- 무엇을 생각했나?} &
\textbf{\footnotesize\color{indigoc} Wonder --- 무엇이 궁금한가?} \\[4pt]
\end{tabularx}
\noindent
\begin{minipage}[t]{0.32\linewidth}\writespace{3}\end{minipage}\hfill
\begin{minipage}[t]{0.32\linewidth}\writespace{3}\end{minipage}\hfill
\begin{minipage}[t]{0.32\linewidth}\writespace{3}\end{minipage}

\section{마무리 성찰}

\begin{sheetbox}
\begin{minipage}[t]{0.48\linewidth}
{\footnotesize\color{inksoft} 오늘 배운 것을 한 문장으로}
\writespace{2}
\end{minipage}\hfill
\begin{minipage}[t]{0.48\linewidth}
{\footnotesize\color{inksoft} 감사 일기 / 유념 한 줄}
\writespace{2}
\end{minipage}
\end{sheetbox}

\end{document}
